\section{总结与延伸}

\subsection{核心启示}

Anthropic 在这篇工程博客中传递的最重要信息是：为 Agent 构建高效工具需要从传统确定性软件开发范式转向面向非确定性系统的新范式。

通过评估驱动的迭代过程，Anthropic 团队识别出了成功工具的一致模式：

\begin{enumerate}[leftmargin=1.5em]
    \item 有效的工具被\textbf{有意识地、清晰地定义}
    \item 它们\textbf{审慎地使用 Agent 上下文}
    \item 它们可以在\textbf{多样化的工作流中灵活组合}
    \item 它们使 Agent 能够\textbf{直觉地解决现实世界任务}
\end{enumerate}

\subsection{展望未来}

文章预期 Agent 与世界交互的具体机制将不断演化——从 MCP 协议的更新到底层 LLM 自身的升级。但评估驱动的工具优化方法论将保持其价值：随着 Agent 变得更加强大，与之配合的工具也将随之进化。

\subsection{拓展阅读}

\begin{itemize}[leftmargin=1.5em]
    \item Anthropic Developer Guide: Tool Use Best Practices\\ \url{https://docs.anthropic.com/en/docs/build-with-claude/tool-use}
    \item Model Context Protocol (MCP) 官方文档\\ \url{https://modelcontextprotocol.io/}
    \item Anthropic Tool Evaluation Cookbook\\ \url{https://github.com/anthropics/anthropic-cookbook}
    \item SWE-bench Verified 评测\\ \url{https://www.swebench.com/}
    \item Claude Code 文档\\ \url{https://docs.anthropic.com/en/docs/claude-code}
\end{itemize}

%% --- Fin del contenido --- %%

\end{document}
